% ------------------------------------------------------------
% Addendum: Euclidean and Non-Euclidean Geometry in A6/A7
% Author: Borealis S. Hedling
% Altitude: A6 → A7
% Purpose: Formalize geometric distinctions required for dual and non-dual
%          operator behavior, collapse-surface geometry, and divergence membranes.
% ------------------------------------------------------------

\section{Geometric Foundations of the Altitude Model}

This addendum formalizes the geometric distinctions between the A6 Euclidean
regime and the A7 non-Euclidean boundary regime. These geometric properties are
structurally required for dual algebra (A6), non-dual algebra (A7), collapse-
surface behavior, divergence invariants, and non-dual projection operators.

% ------------------------------------------------------------
\subsection{A6 Euclidean Geometry}

The A6 manifold $M$ is Euclidean:
\[
K(p) = 0 \quad \forall p \in M,
\]
where $K$ denotes sectional curvature.

Euclidean properties preserved by A6 operators include:
\begin{itemize}
  \item geodesic linearity,
  \item adjacency preservation,
  \item bounded holonomy,
  \item spectral continuity,
  \item modal continuity.
\end{itemize}

Holonomy drift remains bounded:
\[
H(\gamma) \le \epsilon_{\mathrm{holo}}.
\]

Adjacency is preserved under all A6 operators:
\[
A(p,q) = A(O(p), O(q)).
\]

Spectral continuity holds:
\[
|\delta(O(p))| \le \epsilon_{\mathrm{spec}}.
\]

% ------------------------------------------------------------
\subsection{A7 Non-Euclidean Geometry}

The A7 boundary manifold is the divergence membrane:
\[
M = \partial S.
\]

This hypersurface is non-Euclidean:
\[
K(p) > 0 \quad \forall p \in \partial S.
\]

Non-Euclidean properties introduced at A7 include:
\begin{itemize}
  \item curvature accumulation,
  \item geodesic divergence,
  \item adjacency fracture,
  \item spectral tearing,
  \item holonomy accumulation.
\end{itemize}

Holonomy increases under A7 operators:
\[
H(D(\gamma)) = H(\gamma) + \Delta_{\mathrm{holo}}.
\]

Adjacency fractures:
\[
A(p,q) \ne A(D(p), D(q)).
\]

Spectral tearing occurs:
\[
\mathrm{Spec}(D(p)) \not\approx \mathrm{Spec}(p).
\]

% ------------------------------------------------------------
\subsection{Altitude Boundary Geometry}

The altitude boundary between A6 and A7 is the divergence membrane:
\[
M = \partial S.
\]

This boundary separates:
\begin{itemize}
  \item the Euclidean interior ($A6$),
  \item the non-Euclidean boundary ($A7$).
\end{itemize}

Operator domains respect this separation:
\[
O_{\mathrm{A6}} : S \to S,
\qquad
D_{\mathrm{A7}} : \partial S \to \partial S.
\]

% ------------------------------------------------------------
\subsection{Collapse-Surface Geometry}

Divergence collapse surfaces are non-Euclidean hypersurfaces:
\[
\Sigma_{\mathrm{div}} = \{ p \mid I_{\mathrm{div}}(p) = c \},
\]
with curvature:
\[
K(\Sigma_{\mathrm{div}}) > 0.
\]

These surfaces reconcile:
\begin{itemize}
  \item spectral tears,
  \item modal discontinuities,
  \item holonomy fractures,
  \item curvature singularities.
\end{itemize}

% ------------------------------------------------------------
\subsection{Geometry of A7 Non-Dual Projection Operators}

Non-dual projection operators incorporate non-Euclidean geometry:
\[
\Pi_{\mathrm{ND}} : M \to \Sigma.
\]

Curvature increases under projection:
\[
K(\Pi_{\mathrm{ND}}(p)) > K(p).
\]

Holonomy breaks:
\[
H(\Pi_{\mathrm{ND}}(\gamma)) = H(\gamma) + \Delta_{\mathrm{holo}}.
\]

Geodesics fracture:
\[
\gamma \mapsto \gamma' \quad \text{with discontinuous tangent}.
\]

Spectral tearing is preserved:
\[
\mathrm{Spec}(\Pi_{\mathrm{ND}}(p)) \not\approx \mathrm{Spec}(p).
\]

% ------------------------------------------------------------
\section{Addendum Summary}

A6 operates on a Euclidean manifold with continuous, monotone, adjacency-
preserving behavior. A7 operates on a non-Euclidean boundary manifold where
curvature, holonomy, adjacency, and spectral properties diverge. This geometric
distinction is essential for maintaining altitude separation, operator
consistency, and non-dual algebraic behavior.

% ------------------------------------------------------------
\section*{Provenance}

This addendum is part of the Unified Altitude Model and formalizes geometric
requirements for A6 dual algebra and A7 non-dual algebra. All content is
simulation-safe, non-regulatory, and non-dual-use in the export-control sense.
